% ---------------------------------------------------------------- Appendix A
\section{Assumptions}\label{app:A}

\begin{table}[htbp]\centering\footnotesize
\caption{Research question and model assumptions}\label{tab:A1}
\begin{tabularx}{\linewidth}{@{}>{\hsize=0.78\hsize}L>{\hsize=1.02\hsize}L>{\hsize=1.20\hsize}L@{}}
\toprule
Item & Value & Remark \\
\midrule
Research question & Does the corporate tax turn debt into a free lunch? & On trial: no corporate tax; distress costs added only to bound debt \\
Benchmark & $\VL = \VU$ & No taxes, no distress costs, equal borrowing terms (Lecture~1, p.~49) \\
EBIT, strong / weak state & 120 / 60 & Two-state firm of Lecture~1 (pp.~9, 47) \\
State probabilities & 0.5 / 0.5 & Symmetric states \\
Promised payment in benchmark example & 80 & Section~\ref{sec:notax} only \\
Horizon & One period; liquidation at date 1 & Multi-period version in Section~\ref{sec:multi} \\
Discount rate & 10\% & Risk-free rate; risk-neutral probabilities \\
Debt at the starting point & 40 & Base case of the investigation \\
Borrowing rate & 10\%, independent of leverage & Market-priced version in Section~\ref{sec:mkt} \\
Promised repayment & $44 = 40 \times 1.10$ & \\
Corporate tax rate $\tauc$ & 20\% & 17\% and 30\% in Section~\ref{sec:robust} \\
Interest / principal deductible & Yes / No & Source of the shield \\
Loss carry-forward & None & Discussed in Sections~\ref{sec:usable} and~\ref{sec:sg} \\
Personal taxes / trading costs & None / none & Personal taxes in Section~\ref{sec:personal} \\
Distress cost $c$ & 20\% of after-tax cash in default (of $\VU$ in Sections~\ref{sec:multi}--\ref{sec:apps}) & 5\% and 40\% in Section~\ref{sec:robust} \\
Continuous cash flow & Lognormal, mean 90, s.d.\ 25 & 200,000 draws; s.d.\ 30 as a check \\
\bottomrule
\end{tabularx}
\tablenote{\textit{Notes:} All values are illustrative rather than data on a real company, except where Section~\ref{sec:apps} cites a source.}
\end{table}
